\section{Code as Modality：为什么 BLEU 不够}

代码既是 token sequence，又有可执行语义。两个字符串差异很大可以实现相同函数，字符串相似也可能因一个符号产生严重错误。因此 code generation 的评价应尽量执行 specification，而不是只比较 reference overlap。

\subsection{Code 真的只是另一种 modality 吗}

本节回到可自动验证的输出。Code 可以沿用 token generation recipe，但 parser、runtime 与 tests 让“正确”比自然语言相似度更可操作，也更严格。前面的文本和图像多依赖人工或感知评价；代码首次让候选分布能够进入自动执行、失败分类和搜索闭环。

\lecturefigure{slide-29-code-modality.jpg}{Code 可以用语言模型生成，但其正确性由执行语义约束。}{官方视频 32:20。}

\begin{knowledgebox}{读图：相同 recipe，新增 verifier}
Code tokenization 和 autoregressive objective 可直接复用 GPT，但 code 拥有 parser、type checker、tests 与 runtime。它们提供比自然语言更强的自动反馈。与此同时，程序还有 security、resource、dependency 和 specification ambiguity；通过有限 tests 只是证据之一。
\end{knowledgebox}

\teachervoice{讲者问“isn't code just another modality?”随后立刻转向 HumanEval。课堂提示：模型接口可以统一，evaluation 不能偷懒；代码的优势是有 executable verifier，责任也更高。}

\subsection{HumanEval：问题、签名、docstring 与 tests}

\term{HumanEval} 是一组手写 Python function problems。每题给 function signature 与 docstring，模型生成 body；隐藏 \term{unit tests}（单元测试）检查多个输入。\term{functional correctness}（功能正确性）指程序在 specification cases 上产生正确行为，而不是与 reference code 文本相似。

\lecturefigure{slide-30-humaneval-dataset.jpg}{HumanEval 用 function specification 和 hidden tests 评价生成代码。}{官方视频 33:33；Chen et al. 2021。}

\begin{knowledgebox}{读图：一题包含哪些契约}
Signature 规定接口，docstring 给自然语言 specification，examples 可能暗示边界，tests 定义可观察正确性。若 tests 不覆盖异常输入、复杂度或副作用，错误程序仍可能通过。Benchmark quality 取决于 specification clarity 与 test coverage。
\end{knowledgebox}

\lecturefigure{slide-31-humaneval-example.jpg}{HumanEval example 展示 prompt 与候选 implementation 的关系。}{官方视频 35:15。}

\begin{knowledgebox}{读图：先人工推导，再运行 tests}
阅读函数名、参数与 docstring，列出 normal、boundary、empty、duplicate 和 invalid cases；再检查候选代码。模型可能写出语法正确且看似合理的实现，却在 off-by-one、mutation 或类型边界失败。执行是必要验证，但 secure sandbox 与 timeout 也不可缺少。
\end{knowledgebox}

\begin{lstlisting}[caption={HumanEval-style functional check 的简化形式}]
def candidate(values):
    # generated implementation
    ...

assert candidate([1, 2, 3]) == expected_normal
assert candidate([]) == expected_empty
assert candidate(boundary_case) == expected_boundary
\end{lstlisting}

\subsection{Pass@k：分布中至少一个正确样本}

若从模型采样 $k$ 个独立候选，\term{pass@k} 表示至少一个通过 tests 的概率。对每题采样 $n$ 个程序，其中 $c$ 个正确，Codex 论文使用估计量
\[
\widehat{\operatorname{pass@}k}
=1-\frac{\binom{n-c}{k}}{\binom{n}{k}},
\]
当 $n-c<k$ 时结果为 1。分子表示从错误样本中选出 $k$ 个的组合数，分母是所有 $k$ 子集；两者比值是“选出的全部错误”概率。

\lecturefigure{slide-32-pass-at-k.jpg}{Pass@k 评价多次采样中至少找到一个功能正确程序的机会。}{官方视频 36:00；Chen et al. 2021。}

\begin{knowledgebox}{读图：pass@1 与 pass@100 回答不同产品问题}
Pass@1 接近用户只看一次 completion 的体验；large-k 更接近系统可生成许多候选并用 tests/reranker 选择。高 pass@100、低 pass@1 表示正确程序存在于分布，但默认 decoding 很少命中。部署价值取决于是否真有 verifier、采样成本和选择延迟。
\end{knowledgebox}

\begin{warningbox}{Pass@k 不是“模型会写软件”的比例}
HumanEval 题目短、specification 清晰、tests 自动化。真实软件包含跨文件依赖、模糊需求、security、performance、maintenance 与 review。Benchmark pass@k 是窄能力证据，不是 autonomous software-engineering reliability。
\end{warningbox}

\subsection{Codex training snapshot}

课堂 slide 记录当时 Codex 数据为从 5400 万 repositories 收集的 159 GB code，并从不同规模 GPT-3 models fine-tune；tokenizer 额外支持 spaces。这里是论文时期历史描述，不代表后来产品数据或当前政策。

\lecturefigure{slide-33-codex-training-details.jpg}{Codex 在 GPT-3 基础上使用大规模 code corpus 微调并调整 tokenizer。}{官方视频 36:59；Chen et al. 2021。}

\begin{knowledgebox}{读图：数据量之外还有 representation choice}
Repository count 不能直接转成独立、高质量 examples；需要去重、许可、安全过滤与 train-test contamination audit。Spaces 对 Python indentation 有语义，tokenizer 改动可降低序列长度并改善结构学习。Fine-tuning 复用自然语言能力，也把自然语言 prompt 与 code distribution 接起来。
\end{knowledgebox}

\subsection{本章小结}

Code generation 继承 language-model recipe，却因 executable semantics 获得更强 verifier。HumanEval 与 pass@k 把“像代码”升级为“至少一个候选通过 tests”，但仍只覆盖真实软件的一小部分。

